% =====================================================================
\section{Accuracy of computation and verifiability of claims}\label{sec:prec}
Version~5 made the computation three times faster and at the same time established what in it may be believed.

\textbf{Where the speed-up came from.} The computation was limited not by arithmetic but by memory bandwidth: every elementary operation on the array of cells is a separate pass, and a step made about forty-five such passes. The whole per-cell block was rewritten as a single fused loop in which the intermediate quantities live in registers. The gain at the working size is \textbf{3.3-fold in wall time}: a full run of 123 blocks at 2400 sets takes 3 hours 05 minutes instead of nearly four hours previously at a smaller volume. The mechanism is visible in fig.~\ref{fig:v5}d: for the one-operation-at-a-time computation the cost per row doubles between 2000 and 5250 rows --- the working set ceases to fit in cache --- and holds thereafter, whereas for the fused kernel it is flat at all sizes, because its working set is a single row.

\textbf{What proves that this is the same model.} The fast path was not rewritten by hand but \emph{assembled from the source} of the reference by text substitution with a mandatory uniqueness check on every fragment: on any edit to the model the assembly fails rather than diverging silently. The comparison is made at two levels --- all 63 fields of the result at five points (including both sides of the asymmetry; the counting fields agree bitwise, the worst discrepancy of the continuous ones being $9.2\cdot10^{-15}$) --- and separately the whole set of rules, run twice with a diff of all verdicts.

\textbf{What to believe in the time step.} The quantities split into two classes. The continuous means converged with margin: at a step twice as coarse they hold within 0.5\,\%. The counting ones did not: the collapse rate falls monotonically under coarsening of the step in all conditions, and Richardson extrapolation gives a limit \emph{2\,\% above} the value at the working step. The cause is not the moment of threshold crossing: sub-step interpolation refines it but yields only 0.009\,\%. The bias comes from missed sub-step excursions and from the error of the trajectory itself. The conclusion for reading the results: thresholds on counting quantities finer than about 3\,\% may not be believed, and this is recorded as a known limit rather than painted over.

\textbf{The seventh trap, and it has already found its way into the text of this brochure once.} \emph{The mean over a bimodal population.} At the working point the sets fall into two heaps: either broken irreversibly (damaged fraction above 0.9) or untouched (below 0.01), with almost nobody in between. The mean over such a population is the \emph{fraction of broken sets}, not the state of the system, and it may not be read as a state. This is precisely how the claim arose in version~4 that the model does not reach the locking staircase: a mean over the mixture was being compared, whereas in the broken sets the damage stood at unity even then (section~\ref{sec:bifur}). The rule for reading: wherever a damaged fraction appears in the tables, the fractions of sets above 0.9 and below 0.01 must be kept beside it, or the mean will deceive.

\textbf{The fifth and sixth traps.} The fifth is the choice of the measured quantity (section~\ref{sec:burst}). The sixth is an \emph{irrefutable claim of agreement}: an equivalence tolerance exceeding the curve's own range will pass any data. A rule is now obliged to check this too, and on violation returns ``unverifiable'' with a statement of how many sets are needed.

\begin{ineq}
A separate device which paid off at once. Five parts of rules were unverifiable not for want of sets but by an \emph{information limit}: the rank association was measured over eight cells, its error is $1/\sqrt{n-1}=0.38$, so the detectable threshold is 0.94 against the $-0.50$ recorded. There is nothing to subdivide the cells into: the curve is indexed by wakings, and there is about one per day. The way out is to compute the slope for \emph{each set} over its own cells and to test the spread between sets, where $n$ is in the hundreds. The threshold must then be recomputed: the curve of a single set is noisier than the pooled one, and the association is attenuated by the classical factor
\[\rho_{\text{set}}=\rho_{\text{pooled}}\big/\sqrt{1+\lambda},\qquad \lambda=\frac{\text{variance of the cell noise}}{\text{variance of the signal}},\]
with $\lambda$ estimated from the same data \emph{before} we looked at the slopes. On synthetic data with a known answer the formula is accurate to 1--5\,\%. In the event $\lambda=33.8$, the threshold $-0.50$ was recomputed to $-0.085$, and the MDE fell from 0.94 to 0.076 --- twelvefold --- and four parts out of five closed on data already computed.
\end{ineq}

\textbf{Two rules were rebuilt along the way.} The rule on the composition of an episode lived on a \emph{second} integrator, a copy of the step inside the rule itself, with its own self-comparison against the reference. The comparison failed as soon as interpolation entered the reference: the discrepancy came out as exactly one step per episode boundary. The self-comparison worked as intended, but only one conclusion follows from it --- a duplicate diverges at every edit to the model. It has been deleted, the episode records are written by the reference itself, and there is nothing left to diverge. In passing the sample grew eightfold: 438 sets and 5323 episodes instead of 62 and 788. The second is the labelling of wakings: they were assigned to a cell by the \emph{maximum} damage during that same waking, so the label was partly created by the length being measured. The reverse causality is removed by labelling on the state at the \emph{beginning} of the waking.
